\documentclass[conference]{IEEEtran}
\IEEEoverridecommandlockouts

\usepackage{cite}
\usepackage{amsmath,amssymb,amsfonts}
\usepackage{graphicx}
\usepackage{textcomp}
\usepackage{xcolor}
\usepackage{booktabs}
\usepackage[T1]{fontenc}
\usepackage[hidelinks]{hyperref}

\definecolor{outlinegray}{gray}{0.35}
\newenvironment{outline}[1]
  {\par\medskip\noindent\color{outlinegray}\small
   \textbf{Outline.} \textit{#1}\begin{itemize}\setlength{\itemsep}{2pt}}
  {\end{itemize}\par\medskip}

\begin{document}

\title{MA-X3D: A Lightweight Motion-Attentive 3D CNN\\
for Violence Detection in Surveillance Video}

\author{\IEEEauthorblockN{Nguyen Huu Hoang}
\IEEEauthorblockA{\textit{School of Information and Communications Technology} \\
\textit{Hanoi University of Science and Technology}\\
Hanoi, Vietnam \\
hoang.nh225972@sis.hust.edu.vn}
\and
\IEEEauthorblockN{Ngo Thanh Trung}
\IEEEauthorblockA{\textit{School of Information and Communications Technology} \\
\textit{Hanoi University of Science and Technology}\\
Hanoi, Vietnam}
}

\maketitle

\begin{abstract}
Surveillance video is now produced far faster than operators can watch
it, which motivates automatic detection of violent events. On the
RWF-2000 benchmark, the highest reported accuracies come from
transformer and convolution--transformer backbones with tens to
hundreds of millions of parameters. Models of that size are expensive to
run on every camera stream of an installation. This paper keeps a
compact 3D convolutional backbone, X3D-M, and adds two small modules to
it. The first, Motion Attention, computes a multi-scale frame-difference
map, converts it into four learnable temporal modes, and uses them to
gate the intermediate features towards regions with movement. The second
is a wide-kernel convolution that enlarges the spatial support of the
two earliest residual stages from $3\times3$ to $5\times5$; the added
outer ring of weights is trained in its own parameter group at a higher
learning rate and is fused into a single kernel after training. A
person-centred crop prepares the input for the static-camera setting.
Both modules are initialised as an identity, so training starts from the
unmodified pre-trained backbone. On RWF-2000 the model reaches 90.5\%
accuracy and 0.91 F1 on the violent class, 5.75 points above the X3D-M
baseline, using 3.02M parameters and 5.16 GFLOPs per clip. It reaches
97.0\% on the Hockey Fight dataset. An ablation reports the contribution
of each component.
\end{abstract}

\begin{IEEEkeywords}
violence detection, video action recognition, surveillance video,
efficient 3D CNN, motion attention, structural reparameterisation
\end{IEEEkeywords}

\section{Introduction}
\label{sec:intro}

The number of closed-circuit television (CCTV) cameras in public spaces
has grown faster than the capacity to monitor them. By the end of 2019
an estimated 770 million surveillance cameras were in
operation~\cite{cnbc2019}, and a large metropolitan network can produce
hundreds of terabytes of footage per day. Most installations still rely
on operators watching live feeds. Vigilance studies show that this is
unreliable over long shifts: the ability to detect infrequent signals in
a continuous monitoring task starts to fall within the first twenty
minutes and keeps falling afterwards~\cite{mackworth1948}. An automatic
detector does not remove the operator from the loop. It flags candidate
events so that attention can be spent on those rather than on all
streams at once.

Violent events are hard to detect from surveillance footage. They
usually last a few seconds, involve fast and irregular body movement,
and may cover only a small part of a wide-angle frame. A detector also
has to work across many camera viewpoints, resolutions and lighting
conditions, and has to separate violence from non-violent activity that
looks similar, such as animated conversation or horseplay. Work on the
problem has followed video action recognition more generally: from
hand-crafted motion descriptors~\cite{idt,vif,hockey}, to two-stream
networks that consume pre-computed optical flow~\cite{twostream}, to 3D
convolutional networks that learn spatio-temporal features from raw
frames~\cite{i3d,slowfast,x3d}, and more recently to video
transformers~\cite{timesformer,videoswin}.

Cost is a third requirement and it is often left out of the comparison.
An installation may carry tens or hundreds of cameras, so a detector is
only usable if it can run on every stream at the same time. Recent
results on RWF-2000~\cite{rwf2000} have moved away from that constraint.
The highest published accuracy is 94.00\%, obtained by
CUE-Net~\cite{cuenet} with 354M parameters and 5826 GFLOPs per clip.
Masked-autoencoder video transformers report 91.8\% to 92.4\% with 67M
to 305M parameters~\cite{videomae,mixvit}. Accuracy on this benchmark
has therefore been obtained at one to two orders of magnitude more
computation than a per-stream detector can afford. A separate line of
work has aimed for competitive accuracy inside a small
budget~\cite{rwf2000,sepconvlstm,mdcn,biltma}, and this paper follows
that line.

Two gaps remain in how compact 3D backbones are used for violence
detection. First, motion is only implicit. A 3D backbone has to infer
movement from appearance differences between frames and is given no
explicit motion signal, although the frame difference is available at
the cost of a subtraction and was shown by
SepConvLSTM~\cite{sepconvlstm} to work as a substitute for optical flow
on this benchmark. Existing work that uses this cue passes it through a
second input stream, which requires a second backbone and roughly
doubles the cost. Second, the early receptive field is narrow. The
depthwise kernels of X3D have a fixed $3\times3$ spatial support at
every stage, which limits the earliest stages. Those stages run at the
highest spatial resolution and are the ones best placed to resolve
small, fast displacements.

We propose MA-X3D, which addresses both gaps while staying in the
lightweight range. The X3D-M backbone~\cite{x3d} is left unchanged and
two small modules are attached to it. Both are initialised as an
identity, so the model at epoch zero is functionally the pre-trained
backbone and the new parameters are learned as departures from it. This
matters on RWF-2000, where 1600 training clips are too few to learn a
spatio-temporal representation from scratch and randomly initialised
additions would disturb the features transferred from Kinetics
pre-training. The Motion Attention module converts a multi-scale frame
difference into four learnable temporal modes and uses a near-zero gate
to modulate the Res4 features multiplicatively. The motion cue therefore
re-weights features the backbone has already computed, and it does so
inside one backbone instead of two. The wide-kernel convolution enlarges
the depthwise kernels of Res2 and Res3 by structural
reparameterisation~\cite{repvgg,acnet}: a zero-initialised outer ring is
added around the pre-trained centre and is placed in a parameter group
whose learning rate is fifty times that of the backbone. Without this
separation the ring stays close to its initialisation. After training
the two terms fuse into a single $5\times5$ kernel, so no training-time
structure remains at inference.

The contributions of this paper are the following.
\begin{itemize}
  \item A Motion Attention module that adds an explicit
        frame-difference cue to a single 3D backbone through learnable
        temporal modes and an identity-initialised multiplicative gate.
        It adds about 2.5K parameters and no second branch.
  \item A wide-kernel reparameterised convolution that enlarges the
        receptive field of the earliest stages and leaves no structural
        overhead at inference, together with the learning-rate scheme
        that makes the added ring trainable. Naive dense inflation of
        the same kernel gains 2.25 accuracy points less and leaves the
        ring close to zero.
  \item An evaluation on the official RWF-2000 split, reaching 90.5\%
        accuracy and 0.91 F1 on the violent class with 3.02M parameters
        and 5.16 GFLOPs, an ablation for each component, a Grad-CAM and
        gate analysis, and a cross-domain result on Hockey Fight
        (97.0\%).
\end{itemize}

Section~\ref{sec:related} reviews related work,
Section~\ref{sec:method} describes MA-X3D,
Section~\ref{sec:exp} reports the experiments, and
Section~\ref{sec:conclusion} concludes.

\section{Related Work}
\label{sec:related}

\subsection{Violence detection}

Early methods described violent action with hand-crafted features.
Wang and Schmid~\cite{idt} encoded motion along densely sampled
trajectories, Hassner~et~al.~\cite{vif} summarised how optical-flow
magnitudes change over time, and Bermejo~et~al.~\cite{hockey} proposed
acceleration-based descriptors alongside the Hockey Fight dataset.
These representations are cheap and interpretable, but their accuracy
falls when the camera angle, scene type, or lighting differs from the
training conditions.

Two-stream networks~\cite{twostream} replaced the hand-crafted step
with learned features. One stream processes RGB frames and the other
processes pre-computed optical flow, which provides a strong motion
signal. The cost is that dense optical flow is expensive to compute,
often more than the recognition network itself, and this limits
real-time use.

Three-dimensional convolutional networks removed the optical flow
dependency. I3D~\cite{i3d} inflates ImageNet-pretrained 2D filters
into the temporal dimension, SlowFast~\cite{slowfast} runs two
pathways at different frame rates to capture slow semantics and fast
motion separately, and X3D~\cite{x3d} finds a compact architecture
through progressive expansion of six dimensions. These models learn
spatio-temporal features directly from RGB.

Video transformers~\cite{timesformer,videoswin} extended the
self-attention mechanism to space and time and raised accuracy further.
On large benchmarks they outperform 3D convolutional networks, but they
depend on large-scale pre-training. On RWF-2000, which provides only
1600 training clips, that dependence is a practical problem.

Table~\ref{tab:sota} places the methods discussed next to each other
on the RWF-2000 benchmark~\cite{rwf2000}. ConvLSTM~\cite{convlstm}
reaches 77.00\% and C3D~\cite{c3d} 82.75\%. The Flow Gated
Network~\cite{rwf2000}, which uses an optical-flow branch as a learned
gate over the RGB branch, reaches 87.25\%.
SepConvLSTM~\cite{sepconvlstm} processes background-suppressed frames
and frame differences through two MobileNetV2 backbones fused by a
separable convolutional LSTM, reaching 89.75\%; the frame-difference
cue is effective, but the two-backbone construction roughly doubles the
computational cost. A multi-dimensional convolutional network with two
streams~\cite{mdcn} reaches 89.70\% below one million parameters, and
a 2D EfficientNet-B0 with bidirectional motion attention and temporal
shift~\cite{biltma} reaches 90.25\% at 4.48M parameters. At the other
end, VideoMAE vision transformers reach 91.8\% to 92.4\% with 67M to
305M parameters~\cite{videomae,mixvit}, and CUE-Net~\cite{cuenet}
reaches 94.00\% using person-centred cropping, an enhanced UniformerV2
backbone, and a modified efficient additive attention module at 354M
parameters and 5826 GFLOPs per clip.

\subsection{Efficient video backbones}

X3D~\cite{x3d} generates a family of models by progressively enlarging
a small 2D network one axis at a time along six axes: temporal
duration, frame rate, spatial resolution, depth, channel width, and
bottleneck width. At each step the axis with the best accuracy-to-cost
ratio is kept. The medium variant, X3D-M, attains 76.0\% top-1
accuracy on Kinetics-400 with 3.8M parameters and 6.2 GFLOPs per view.
This is on par with SlowFast R50 (75.6\%, 34.4M) and Non-local R50
(76.5\%, 35.3M) while using about nine times fewer parameters. X3D-S
and X3D-M share an identical parameter count and differ only in input
resolution; parameters grow only from X3D-XL onward. X3D-M is
therefore the most accurate configuration within the minimal parameter
budget, which is why it is used as the backbone here.

\subsection{Motion cue, attention, and reparameterisation}

Under a static camera the background does not move, so the absolute
difference between consecutive frames isolates the movement of the
people in the scene. This is cheaper than optical flow and does not
require camera-motion compensation. SepConvLSTM~\cite{sepconvlstm}
showed that it is an effective motion cue for violence detection.

Multiplicative gating, used in Squeeze-and-Excitation
networks~\cite{se} and CBAM~\cite{cbam}, re-weights feature maps with
a learned signal. A gate $g$ produces the modulated output
$x \odot (1+g)$. When the last layer of the gate is initialised to
near-zero weights, $g \approx 0$ and the module starts as an identity,
which lets new parameters be added to a pre-trained backbone without
disturbing the transferred features.

Structural reparameterisation exploits the linearity of convolution.
Because a sum of convolutions over the same input is equivalent to a
single convolution with the summed kernel, extra branches used during
training can be fused into one kernel at inference. RepVGG~\cite{repvgg}
trains a multi-branch block and fuses it into a plain convolution for
deployment; ACNet~\cite{acnet} adds asymmetric kernels that merge back
into the original square kernel. Both methods fuse back to the original
kernel shape. In the wide-kernel construction of
Section~\ref{sec:method}, the added ring fuses into a larger kernel,
so the extra spatial support is retained.

\section{Proposed Method}
\label{sec:method}

\subsection{Overview}

The task is clip-level binary classification. A trimmed surveillance
clip is represented by an RGB tensor $x \in \mathbb{R}^{3 \times T
\times H \times W}$ and a single-channel motion map $m \in
\mathbb{R}^{1 \times T \times H \times W}$, with $T=16$ and
$H=W=224$. The model $f_\theta$ maps these to two scores and the
prediction is $\hat{y} = \arg\max f_\theta(x, m)$.

Fig.~\ref{fig:pipeline} shows the pipeline. An offline stage builds
a cached dataset: frames are extracted, a person-centred crop is
computed and applied, and contrast is enhanced. At training and
inference time, 16 frames are sampled from the stored clip, the motion
map is computed from the augmented frames, and the RGB tensor and
motion map are fed to the backbone.

Three principles govern the design. \textbf{P1~Lightweight}: no second
branch; the modules reuse depthwise convolutions and add about 36K
parameters in total. \textbf{P2~Identity initialisation}: each module
produces an identity transformation at epoch zero, so the augmented
model starts from the unmodified pre-trained backbone.
\textbf{P3~No inference-time overhead}: any extra structure used during
training is fused away before deployment.

\begin{figure}[t]
  \centering
  % \includegraphics[width=\linewidth]{fig/pipeline.pdf}
  \caption{Overall pipeline of MA-X3D. Offline preprocessing builds the
  cached dataset; the backbone applies wide-kernel convolutions in Res2
  and Res3 and the Motion Attention module after Res4.}
  \label{fig:pipeline}
\end{figure}

\subsection{Person-centred region of interest}

A five-second window centred on the clip is sampled at 64 uniformly
spaced frames. YOLOv8n~\cite{yolov8} is run on 12 of them to detect
people, and the centres of all detected boxes are clustered with
DBSCAN~\cite{dbscan} using a neighbourhood radius of
$0.15 \max(H, W)$ and $\mathrm{minPts}=2$. The bounding box of the
most populated cluster, enlarged by a 20\% margin on each side, is
used as the crop. The full frame is kept when no person is detected.

Clustering rather than taking the union of all detected boxes is a
deliberate choice. Bystanders standing apart from the action would
extend the union box to include them, diluting the motion signal the
crop is meant to concentrate.

The crop is a single static box applied to every frame of the clip.
Because the camera is fixed and the box does not move, the background
stays registered across frames. The frame difference computed in the
next section therefore reflects genuine person motion rather than
apparent motion caused by a shifting window. A per-frame tracking crop
would break this property.

Each cropped frame is denoised with a bilateral filter and
contrast-enhanced with CLAHE applied to the luminance channel (clip
limit 2.0, $8 \times 8$ tiles), then resized to $224 \times 224$. At
load time, augmentations including brightness and saturation shifts,
horizontal flip, random crop, rotation, and Gaussian blur are applied
identically to all frames of a clip. Applying the same transform to
every frame preserves the inter-frame registration on which the motion
map depends.

\subsection{Motion Attention}
\label{sec:ma}

Let $d^{(s)}_t = \frac{1}{3} \sum_{c=1}^{3} |x_{c,t} - x_{c,t-s}|$
be the channel-averaged absolute frame difference at step $s$ and time
$t$. Using two temporal scales $s \in \{1, 2\}$ captures both fast and
slower motion. The single-channel motion map is
\begin{equation}
  m_t = \min\!\left( \max\!\left(d^{(1)}_t,\, d^{(2)}_t\right),\; 0.2 \right) \times 5 \,.
  \label{eq:motion}
\end{equation}
The threshold 0.2 saturates large values produced by illumination
flicker or compression artefacts while mapping typical inter-frame
differences to approximately $[0, 1]$. Taking the element-wise maximum
rather than concatenating keeps the map single-channel; a location is
marked as moving if either scale detects motion.

A temporal convolution with kernel size $(3, 1, 1)$ maps the
single-channel map to $N=4$ channels:
$M = \mathrm{TConv}(m) \in \mathbb{R}^{4 \times T \times H \times W}$.
Because the spatial extent is one, this convolution mixes information
only along time, so each output channel can represent a different
temporal pattern. The first mode is initialised as a temporal impulse
on the current frame, reproducing the raw motion map, and the remaining
three from $\mathcal{N}(0, 0.05^2)$.

The four modes are resized to the Res4 feature resolution by trilinear
interpolation, giving $\tilde{M} \in \mathbb{R}^{4 \times T' \times 14
\times 14}$. A gating network $g$ maps $\tilde{M}$ to a modulation:
$1 \times 1 \times 1$ convolution to 24 channels, ReLU, $1 \times 1
\times 1$ convolution to 96 channels, then tanh. The modulated Res4
feature map $f$ is
\begin{equation}
  y = f \odot \bigl(1 + g(\tilde{M})\bigr) \,.
  \label{eq:gate}
\end{equation}
Since $g \in [-1, 1]$, a position is amplified where $g > 0$,
suppressed where $g < 0$, and unchanged where $g = 0$. The final
convolution of the gate is initialised from $\mathcal{N}(0, 0.01^2)$
with zero bias, so $g \approx 0$ and $y \approx f$ at the start of
training, satisfying P2.

The module is placed after Res4 because at $14 \times 14$ the feature
map still resolves spatial position while the features are already
semantic. Res2 and Res3 operate on lower-level texture, and Res5 at
$7 \times 7$ is too coarse to localise where motion occurs. The
residual is multiplicative rather than additive: multiplication
re-weights features the backbone has already computed, while addition
would inject new content that competes with the pre-trained
representation. The multiplier $1 + g$ is bounded to $[0, 2]$, so the
motion signal cannot overwhelm the appearance features.

\begin{figure}[t]
  \centering
  % \includegraphics[width=\linewidth]{fig/ma.pdf}
  \caption{The Motion Attention module. The motion map is converted to
  four temporal modes, interpolated to the Res4 resolution, and gated
  to modulate the features as $f \odot (1+g)$.}
  \label{fig:ma}
\end{figure}

\subsection{Wide-kernel reparameterised convolution}
\label{sec:wk}

The depthwise convolutions in Res2 and Res3 use a $3 \times 3 \times
3$ kernel and operate at the highest spatial resolutions in the network,
$56 \times 56$ and $28 \times 28$. Widening their spatial support from
$3 \times 3$ to $5 \times 5$ increases the receptive field of these
stages, which is where fine and fast spatial displacements are best
resolved.

A naive approach inflates each kernel to a dense $5 \times 5$ kernel
with the outer ring initialised to zero. This does not work: the ring
is treated as part of the same parameter group as the backbone and
receives only the small backbone learning rate. Under AdamW the
effective step size is controlled by the learning rate rather than the
raw gradient scale, so at $10^{-5}$ the ring barely moves. The ablation
in Section~\ref{sec:exp} confirms this.

The widened convolution is instead expressed as
\begin{equation}
  W_{\text{eff}} = \mathrm{Pad}_{3 \to 5}(W_{\text{base}})
    + M_{\text{ring}} \odot W_\delta \,,
  \label{eq:rep}
\end{equation}
where $W_{\text{base}}$ is the pre-trained $3 \times 3$ kernel padded
to $5 \times 5$, $W_\delta$ is a $5 \times 5$ kernel initialised to
zero, and $M_{\text{ring}}$ is a fixed binary mask that is one on the
16 outer positions of the $5 \times 5$ grid and zero on the central
$3 \times 3$. The mask ensures that the centre of the effective kernel
always equals the pre-trained weights, and the zero initialisation
makes the convolution an identity at the start of training.

Because $W_{\text{base}}$ and $W_\delta$ are separate parameters,
$W_\delta$ can be placed in a high-learning-rate group
($5 \times 10^{-4}$) while $W_{\text{base}}$ remains at $10^{-5}$.
This separation is what makes the outer ring trainable. Since
convolution is linear in its kernel, the two terms in~(\ref{eq:rep})
fuse after training into a single $5 \times 5$ kernel, leaving no
structural overhead at inference. The widening is applied to all eight
depthwise convolutions in Res2 (three blocks) and Res3 (five blocks),
adding 33{,}696 parameters.

Widening the depthwise kernel from $k=3$ to $k=5$ doubles the
per-convolution contribution to the receptive field from $k-1=2$ to
$k-1=4$ pixels. A stage of $L$ blocks therefore grows from $2L$ to
$4L$ pixels: from 6 to 12 across Res2 and from 10 to 20 across Res3.

\begin{figure}[t]
  \centering
  % \includegraphics[width=\linewidth]{fig/wk.pdf}
  \caption{The wide-kernel reparameterised convolution. During training
  the $5 \times 5$ kernel is the frozen $3 \times 3$ centre padded to
  $5 \times 5$ plus a zero-initialised outer ring trained at a higher
  learning rate. At inference the two fuse into a single $5 \times 5$
  kernel.}
  \label{fig:wk}
\end{figure}

\subsection{Two-phase discriminative fine-tuning}

Training combines a large pre-trained backbone with several randomly
initialised modules. The schedule adapts the new components first
before touching the backbone.

In the probe phase (four epochs), the entire backbone is frozen and
only the classification head and the Motion Attention module are
trained at a learning rate of $10^{-3}$ with a one-epoch linear
warm-up followed by cosine annealing. This aligns the new components
with the frozen features before any backbone weight moves.

In the fine-tuning phase, Res2, Res3, and Res5 are unfrozen together
with the head, while Res4 stays frozen. Res4 contains no newly added
parameters, and keeping it frozen reduces the number of trainable
backbone weights. Two parameter groups are used: pre-trained backbone
parameters at $10^{-5}$ and new parameters (head, Motion Attention
module, and outer-ring kernels) at $5 \times 10^{-4}$. The optimiser
is rebuilt at the phase transition to reset the adaptive moment
estimates. The loss is cross-entropy with label smoothing
$\varepsilon = 0.1$.

\section{Experiments}
\label{sec:exp}

\begin{outline}{1100 words. 90.5\% at 3.02M parameters, the highest in the lightweight group}
  \item Datasets, protocol, metrics.
  \item Main results and comparison table.
  \item Ablation, including naive inflation against reparameterisation.
  \item Qualitative analysis and error modes.
  \item Cross-domain result and deployment cost.
  \item Limitations.
\end{outline}

\section{Conclusion}
\label{sec:conclusion}

\begin{outline}{120 words}
  \item The two contributions, the headline numbers, and three
        directions for future work.
\end{outline}

\section*{Open items}

\begin{outline}{delete before submission}
  \item Four figures to export from the thesis.
  \item Missing GFLOPs for two methods in the cost table.
  \item Single run on a single split so far.
\end{outline}

\bibliographystyle{IEEEtran}
\bibliography{refs}

\end{document}